\section{法规监管：欧盟 AI 法案}

\subsection{EU AI Act 概述}

讲者重点介绍了欧盟人工智能法案（EU AI Act）——这是全球第一部全面的 AI 监管法律框架。

\begin{importantbox}{EU AI Act 核心要素}
\begin{itemize}
    \item \textbf{生效时间}：2024 年 8 月正式通过，分阶段实施，给予企业合规过渡期
    \item \textbf{核心原则}：基于 \textbf{Risk Level}（风险等级）的分级监管
    \item \textbf{监管逻辑}：风险越高的 AI 项目，受到的限制越严格
    \item \textbf{最高级别}：某些应用被直接\textbf{禁止}（Banned）
\end{itemize}
\end{importantbox}

\subsection{风险分级与禁止项}

EU AI Act 将 AI 应用分为多个风险等级，并对不同等级施加不同程度的监管要求。其中，以下应用被明确\textbf{禁止}：

\begin{table}[H]
\centering
\caption{EU AI Act 中被禁止的 AI 应用类型}
\begin{tabular}{@{}ll@{}}
\toprule
\textbf{禁止项} & \textbf{说明} \\
\midrule
公共场所生物识别 & 实时面部识别用于公共监控 \\
社会评分系统 & 基于个人/群体特征的社会信用评分 \\
操纵性 AI & 利用人类弱点进行潜意识操纵 \\
预测性执法 & 仅基于画像进行犯罪预测 \\
\bottomrule
\end{tabular}
\end{table}

讲者特别提到了\textbf{社会评分系统}（Social Scoring Systems）——一种给个人或群体分配数字评分的系统，这些评分可能被用于信贷审批、就业筛选等高影响决策。他指出，美国实际上已经存在类似系统，例如用于预测累犯率（Recidivism）的算法，根据个人特征数据预测其再犯可能性。在 EU AI Act 框架下，这类系统将被视为非法。

\begin{warningbox}{累犯率预测系统的偏见问题}
美国的累犯率预测系统（如 COMPAS）已被广泛批评存在种族偏见——它会系统性地高估非裔美国人的再犯概率，同时低估白人的再犯概率。这是``看起来客观的算法''实际放大了社会偏见的典型案例。EU AI Act 明确将此类系统列为禁止项。
\end{warningbox}

\subsection{可解释性要求与技术困境}

EU AI Act 中包含一项讲者认为存在问题的条款：\textbf{Explainability}（可解释性）要求。法案要求 AI 系统必须能够解释其决策过程——例如，自动驾驶汽车为什么停车？为什么转弯？

讲者从两个层面分析了这一要求的困难：

\begin{enumerate}
    \item \textbf{数学层面的解释是可能的}：我们可以精确描述神经网络的每一步计算——激活值乘以权重矩阵、应用激活函数、与偏置比较。但这显然不是法规所要求的解释层次。
    \item \textbf{人类可理解的解释是（理论上）不可能的}：深度学习系统属于复杂系统（Complex Systems）——它们具有涌现性（Emergence）、自组织性（Self-organization）和自适应性（Adaptivity）。对于这类系统，在人类可理解的抽象层次上给出决策原因，可能是\textbf{理论上不可能}的。
\end{enumerate}

\begin{importantbox}{复杂系统与可解释性悖论}
深度学习系统是复杂自适应系统（Complex Adaptive Systems）。研究复杂系统的科学家知道，这类系统的行为是从大量简单组件的交互中涌现出来的，无法简单还原为单个组件的行为。因此，要求``解释为什么 AI 做出了这个决定''，在本质上可能等同于要求``解释为什么天气是这样的''——我们可以跑模拟，但无法给出简洁的因果叙事。
\end{importantbox}

\subsection{本章小结}

EU AI Act 代表了全球 AI 监管的重要里程碑。它采用基于风险等级的分级监管策略，禁止了最高风险的应用（生物识别监控、社会评分），并对高风险应用提出了严格的合规要求。然而，可解释性要求在技术上面临根本性挑战——复杂系统的行为可能无法在人类可理解的层次上被解释。

